% ╔══════════════════════════════════════════════════════════════════════════╗
% ║   ÍNDICE DETALLADO — FÍSICA PREUNIVERSITARIA                             ║
% ║   7 Bloques Temáticos · 45 Temas                                         ║
% ╚══════════════════════════════════════════════════════════════════════════╝

{
\hypersetup{linkcolor=AzulPizarra}
\tableofcontents
}
\newpage

\thispagestyle{fancy}

\begin{center}
  {\Large\sffamily\bfseries\color{AzulPizarra} Índice Temático — Física Preuniversitaria}\\[4pt]
  {\small\sffamily\color{GrisMedio} 7 Bloques · 45 Temas · Mecánica, Estática, Hidrostática, Termodinámica, Electricidad y Óptica}
\end{center}

\vspace{0.5cm}

% ─── Bloque: MECÁNICA ─────────────────────────────────────────────────────
\begin{tcolorbox}[enhanced, colback=AzulPizarra!90, colframe=DoradoAcadémico,
    arc=3pt, boxrule=1.5pt, left=8pt, right=8pt, top=6pt, bottom=6pt,
    fontupper=\sffamily\bfseries\large\color{white}]
  ⬡\quad MECÁNICA
\end{tcolorbox}

% ── Unidad I ────────────────────────────────────────────────────────────
\begin{tcolorbox}[enhanced, colback=AzulPizarra, colframe=DoradoAcadémico,
    arc=2pt, boxrule=0.8pt, left=6pt, right=6pt, top=3pt, bottom=3pt,
    fontupper=\sffamily\bfseries\small\color{white}]
  UNIDAD I \hfill Magnitudes y Movimiento \hfill Temas 1–8
\end{tcolorbox}
\begin{multicols}{2}
  \begin{enumerate}[label=\textcolor{AzulPizarra}{\textbf{\arabic*.}},
      leftmargin=2em, itemsep=1pt]
    \item Magnitudes físicas
    \item Sistema Internacional
    \item Vectores
    \item Movimiento rectilíneo uniforme
    \item Movimiento uniformemente variado
    \item Caída libre
    \item Movimiento parabólico
    \item Movimiento circular
  \end{enumerate}
\end{multicols}

\medskip

% ── Unidad II ───────────────────────────────────────────────────────────
\begin{tcolorbox}[enhanced, colback=AzulMedio, colframe=DoradoAcadémico,
    arc=2pt, boxrule=0.8pt, left=6pt, right=6pt, top=3pt, bottom=3pt,
    fontupper=\sffamily\bfseries\small\color{white}]
  UNIDAD II \hfill Dinámica \hfill Temas 9–12
\end{tcolorbox}
\begin{multicols}{2}
  \begin{enumerate}[label=\textcolor{AzulMedio}{\textbf{\arabic*.}},
      leftmargin=2em, itemsep=1pt, start=9]
    \item Leyes de Newton
    \item Fuerza de rozamiento
    \item Plano inclinado
    \item Dinámica circular
  \end{enumerate}
\end{multicols}

\medskip

% ── Unidad III ──────────────────────────────────────────────────────────
\begin{tcolorbox}[enhanced, colback=AzulPizarra, colframe=DoradoAcadémico,
    arc=2pt, boxrule=0.8pt, left=6pt, right=6pt, top=3pt, bottom=3pt,
    fontupper=\sffamily\bfseries\small\color{white}]
  UNIDAD III \hfill Trabajo y Energía \hfill Temas 13–17
\end{tcolorbox}
\begin{multicols}{2}
  \begin{enumerate}[label=\textcolor{AzulPizarra}{\textbf{\arabic*.}},
      leftmargin=2em, itemsep=1pt, start=13]
    \item Trabajo mecánico
    \item Energía cinética
    \item Energía potencial
    \item Conservación de la energía
    \item Potencia
  \end{enumerate}
\end{multicols}

\medskip

% ── Unidad IV ───────────────────────────────────────────────────────────
\begin{tcolorbox}[enhanced, colback=AzulMedio, colframe=DoradoAcadémico,
    arc=2pt, boxrule=0.8pt, left=6pt, right=6pt, top=3pt, bottom=3pt,
    fontupper=\sffamily\bfseries\small\color{white}]
  UNIDAD IV \hfill Cantidad de Movimiento \hfill Temas 18–20
\end{tcolorbox}
\begin{multicols}{2}
  \begin{enumerate}[label=\textcolor{AzulMedio}{\textbf{\arabic*.}},
      leftmargin=2em, itemsep=1pt, start=18]
    \item Impulso
    \item Cantidad de movimiento
    \item Choques
  \end{enumerate}
\end{multicols}

\vspace{0.3cm}

% ─── Bloques especiales ───────────────────────────────────────────────────
\begin{multicols}{2}

  % ESTÁTICA
  \begin{tcolorbox}[enhanced, colback=AzulPizarra!90, colframe=DoradoAcadémico,
      arc=3pt, boxrule=1pt, left=6pt, right=6pt, top=4pt, bottom=4pt,
      fontupper=\sffamily\bfseries\color{white}]
    ⬡\quad ESTÁTICA — Temas 21–23
  \end{tcolorbox}
  \begin{enumerate}[label=\textcolor{AzulPizarra}{\textbf{\arabic*.}},
      leftmargin=2em, itemsep=1pt, start=21]
    \item Equilibrio de fuerzas
    \item Momento de una fuerza
    \item Centro de gravedad
  \end{enumerate}

  \columnbreak

  % HIDROSTÁTICA
  \begin{tcolorbox}[enhanced, colback=AzulMedio, colframe=DoradoAcadémico,
      arc=3pt, boxrule=1pt, left=6pt, right=6pt, top=4pt, bottom=4pt,
      fontupper=\sffamily\bfseries\color{white}]
    ⬡\quad HIDROSTÁTICA — Temas 24–27
  \end{tcolorbox}
  \begin{enumerate}[label=\textcolor{AzulMedio}{\textbf{\arabic*.}},
      leftmargin=2em, itemsep=1pt, start=24]
    \item Densidad
    \item Presión
    \item Principio de Pascal
    \item Principio de Arquímedes
  \end{enumerate}

\end{multicols}

\vspace{0.3cm}

% TERMODINÁMICA
\begin{tcolorbox}[enhanced, colback=AzulPizarra!90, colframe=DoradoAcadémico,
    arc=3pt, boxrule=1pt, left=8pt, right=8pt, top=5pt, bottom=5pt,
    fontupper=\sffamily\bfseries\color{white}]
  ⬡\quad TERMODINÁMICA — Temas 28–32
\end{tcolorbox}
\begin{multicols}{2}
  \begin{enumerate}[label=\textcolor{AzulPizarra}{\textbf{\arabic*.}},
      leftmargin=2em, itemsep=1pt, start=28]
    \item Temperatura
    \item Dilatación
    \item Calor
    \item Cambios de estado
    \item Gases ideales
  \end{enumerate}
\end{multicols}

\vspace{0.3cm}

% ELECTRICIDAD Y MAGNETISMO
\begin{tcolorbox}[enhanced, colback=AzulMedio, colframe=DoradoAcadémico,
    arc=3pt, boxrule=1pt, left=8pt, right=8pt, top=5pt, bottom=5pt,
    fontupper=\sffamily\bfseries\color{white}]
  ⬡\quad ELECTRICIDAD Y MAGNETISMO — Temas 33–40
\end{tcolorbox}
\begin{multicols}{2}
  \begin{enumerate}[label=\textcolor{AzulMedio}{\textbf{\arabic*.}},
      leftmargin=2em, itemsep=1pt, start=33]
    \item Carga eléctrica
    \item Ley de Coulomb
    \item Campo eléctrico
    \item Potencial eléctrico
    \item Corriente eléctrica
    \item Ley de Ohm
    \item Circuitos eléctricos
    \item Magnetismo
  \end{enumerate}
\end{multicols}

\vspace{0.3cm}

% ÓPTICA
\begin{tcolorbox}[enhanced, colback=AzulPizarra!90, colframe=DoradoAcadémico,
    arc=3pt, boxrule=1pt, left=8pt, right=8pt, top=5pt, bottom=5pt,
    fontupper=\sffamily\bfseries\color{white}]
  ⬡\quad ÓPTICA — Temas 41–45
\end{tcolorbox}
\begin{multicols}{2}
  \begin{enumerate}[label=\textcolor{AzulPizarra}{\textbf{\arabic*.}},
      leftmargin=2em, itemsep=1pt, start=41]
    \item Reflexión
    \item Refracción
    \item Espejos
    \item Lentes
    \item Instrumentos ópticos
  \end{enumerate}
\end{multicols}

\vspace{0.5cm}

\begin{cajatip}
  En los exámenes de admisión (UNI, San Marcos, PUCP), la Mecánica representa
  aproximadamente el 50\% de los problemas de Física. Dominar vectores,
  cinemática y las leyes de Newton es prioritario. Se recomienda no avanzar
  a Dinámica sin dominar el diagrama de cuerpo libre.
\end{cajatip}

\newpage